% answers_short.tex —— 2007 年·普通高等学校招生考试·辽宁卷·理科 · 简短答案速查
% 由 tools/make_short_answers.py 自 answers.tex 的 \gksolhead/\gkanswer 提取，
% 并与 2007/2007_数学真题_辽宁理_答案要点.md 逐题比对后生成，勿手改（改源 MD 后重跑）。

\begin{gknotice}
\textbf{使用说明}：本节只给最终答案便于快速核对，完整推导见解析版。答案为编写者独立推导，\textbf{非官方标准答案}。
\end{gknotice}

\gkpart{选择题}
\begin{center}\begin{tabular}{ccccccccccccc}
\toprule
题号 & 1 & 2 & 3 & 4 & 5 & 6 & 7 & 8 & 9 & 10 & 11 & 12 \\
\midrule
答案 & B & C & D & B & B & A & C & A & D & A & B & C \\
\bottomrule
\end{tabular}\end{center}

\gkpart{填空题}
\begin{itemize}[leftmargin=2.6em]
  \item[\textbf{13.}] $-1$。
  \item[\textbf{14.}] $2$。
  \item[\textbf{15.}] $4\sqrt3\pi$。
  \item[\textbf{16.}] $30$。
\end{itemize}

\gkpart{解答题}
\begin{itemize}[leftmargin=2.6em]
  \item[\textbf{17.}] （1）$[-3,1]$；（2）$\omega=2$，增区间 $\left[-\dfrac{\pi}6+k\pi,\dfrac{\pi}3+k\pi\right]$（$k\in\mathbf Z$）。
  \item[\textbf{18.}] （1）证明见解析；（2）$MA=\dfrac{\sqrt3a}6$，点 $C$ 到平面 $MDE$ 的距离为 $\dfrac a4$。
  \item[\textbf{19.}] （1）$L_1=-\dfrac{q^3}3+144q-10$，$L_2=-\dfrac{q^3}3+81q-10$，$L_3=-\dfrac{q^3}3+50q-10$；（2）$E\xi_q=-\dfrac{q^3}3+100q-10$；（3）$q=10$ 时 $E\xi_q$ 最大，最大值为 $\dfrac{1970}3$。
  \item[\textbf{20.}] （1）$(x-4)^2+y^2=16$；（2）最大值 $-\dfrac{16}9$，最小值 $-8$。
  \item[\textbf{21.}] （1）$t\in(-2,0)\cup(0,2)$，$\lim\limits_{n\to\infty}a_n=\dfrac2{2-t}$；（2）证明见解析。
  \item[\textbf{22.}] （1）（2）（3）证明见解析。
\end{itemize}
